\section{Harmonic representatives and intersection cohomology}\label{sec:hodge}

The domain analysis has produced operators preserving the space of
square integrable harmonic forms on the regular locus. To use those
operators in the algebraic criterion, we need to identify the
cohomology they represent and the action of ordinary classes on it.
We will also identify the integration pairing; in Section
\ref{sec:monodromy} this pairing, together with the Hodge star,
will control the part of monodromy commuting with the LLV algebra.
These compatibilities are established through the local de Rham
complex, so that they refer to the same canonical identification.

\subsection{The local de Rham complex}
Let $\mathcal L^j$ be the sheaf of forms $u$ on $M_{\reg}$ for
which $u$ and its weak differential are square integrable on the
regular part of every compact subset of the open set in question.
Lipschitz partitions of unity make these sheaves fine. Since $M$
is compact, their global complex is $\Dom d_{\max}$.
Quasi-isometric regular-locus diffeomorphisms identify these
complexes of sheaves.

The local calculation has a simple topological interpretation.
A punctured cone remembers the cohomology of its link in every
degree. Extending square integrable forms across the vertex
retains precisely the degrees below the middle threshold. That
truncation, with its natural attaching map, is the local condition
defining middle-perversity intersection cohomology.

We recall explicitly the two local calculations needed. For an
interval factor $t$, write $u=\alpha+dt\wedge\beta$ and choose
$\phi\in C_c^\infty$ with integral one. The maps
\[
 Au=\int\phi(t_0)\alpha(t_0)\,dt_0,\qquad
 Ku(t)=\int\phi(t_0)\int_{t_0}^t\beta(s)\,ds\,dt_0
\]
satisfy $dK+Kd=1-\mathrm{pr}^*A$ and $dA=Ad$.
They are bounded on compact subintervals, and globally on bounded
intervals. Testing against smooth forms proves the identities
for weak derivatives. Iteration removes a cube factor.
For a truncated cone with odd-dimensional link $Z$ and
finite-dimensional $L^2$ cohomology, the cone formula is
\begin{equation}\label{eq:cone-cohomology}
 H^j_{(2)}(C_R(Z))=
 \begin{cases}H^j_{(2)}(Z),&j<(f+1)/2,\\
 0,&j\ge(f+1)/2,
 \end{cases}\qquad f=\dim_{\RR}Z.
\end{equation}
In the first line the map is pullback. This is Cheeger's formula
in the generality of \cite[Proposition~2.4]{bhs92}, applied to
$Z_{\reg}$; regularisation identifies the smooth and weak
de Rham complexes used in the two formulations. In particular,
the equality degree belongs to the vanishing range.

\begin{lemma}\label{lem:local-intersection}
Let $P=\RR^N\times C(Z)$, with $Z$ as in
\eqref{eq:cone-cohomology}, and let
$\iota:P\setminus(\RR^N\times\{0\})\hookrightarrow P$.
At a point of the vertex stratum, the stalk of the $j$th
cohomology sheaf of $\mathcal L^\bullet$ is $H^j_{(2)}(Z)$
for $j\le(f-1)/2$ and is zero above this degree. These
cohomology sheaves are locally constant along the stratum.
The stalk of $\mathcal H^j(R\iota_*\iota^*\mathcal L^\bullet)$
is $H^j_{(2)}(Z)$ in every degree, and the attaching map from
$\mathcal H^j(\mathcal L^\bullet)$ is the identity under these
identifications for $j\le(f-1)/2$.
\end{lemma}
\begin{proof}
Choose shrinking neighbourhoods $Q_\epsilon\times C_\epsilon(Z)$,
where $Q_\epsilon$ is a cube in the smooth factor. Restriction
from the globally $L^2$ complex on such a neighbourhood to the
locally $L^2$ complex, and then to the globally $L^2$ complex
on a smaller neighbourhood, interleaves their direct systems.
They therefore compute the same stalks. The interval homotopy
removes $Q_\epsilon$, and \eqref{eq:cone-cohomology} computes
the remaining cone cohomology. Below the threshold the
identifications all come from pullback of closed forms on $Z$;
they are compatible with restriction and constant in the smooth
coordinate. This also proves local constancy.

After deleting the vertex stratum, the neighbourhood is
$Q_\epsilon\times(0,\epsilon)\times Z$. On compact radial
subintervals its metric is uniformly equivalent to the product
metric. The locally $L^2$ complex is consequently the product
complex there, and the interval homotopy gives $H^j_{(2)}(Z)$
in every degree. Fineness computes the derived direct image by
this complex. In the lower degrees the map from the unpunctured
neighbourhood is induced by the same link pullback, so it is the
identity. Thus the stalk calculation also identifies the
attaching morphism.
\end{proof}

\begin{theorem}\label{thm:T1}
There is a canonical isomorphism
\[
 \Phi:\IH^\bullet(M;\RR)\longrightarrow
 H^\bullet_{(2)}(M_{\reg},g).
\]
The maximal exterior derivative has closed range, and every
cohomology class has a unique harmonic representative. If
$\mathcal A^\bullet$ is the sheaf of locally bounded forms
with locally bounded weak differential, then it resolves
$\RR_M$, and the product
$\mathcal A^\bullet\otimes\mathcal L^\bullet\to
\mathcal L^\bullet$ realises the ordinary-cohomology action
under $\Phi$.
\end{theorem}
\begin{proof}
\emph{Identifying the intersection complex.}
On the regular locus, $\mathcal L^\bullet$ resolves $\RR$.
Along $\Sigma\setminus\Delta$ use the imported chart with
link $Z_1$, and along $\Delta$ the chart with link $Z_2$.
Proposition \ref{prop:models} supplies finite-dimensional
cohomology of the latter. The real codimensions are $8k-4$
and $12k$, so the middle perversities are $4k-3$ and $6k-1$.
These are exactly $(f-1)/2$ for the respective links.
Lemma \ref{lem:local-intersection} verifies the constructibility,
stalk vanishing, and attachment axioms of the intersection
complex \cite[\S3.3, Theorem~3.5]{gm83}. With the convention
that its restriction to the regular locus is in degree zero,
\[
 \mathcal L^\bullet\simeq\IC_M.
\]
The identification restricts to the identity on the constant
sheaf there and is canonical. Fineness gives $\Phi$.

\emph{Closed range and harmonic representatives.}
Finite-dimensional intersection cohomology implies closed
ranges in this Hilbert complex. Here is the functional-analytic
point. For a closed differential $d$ whose algebraic range has
finite codimension in the closed space $\ker d_{\rm next}$,
choose a finite-dimensional algebraic complement $F$ in that
kernel. The bounded map
$\Dom d\oplus F\to\ker d_{\rm next}$, $(u,f)\mapsto du+f$,
from the graph-norm Banach space is surjective. The open mapping
theorem makes the projection to $F$ continuous; its kernel is
$\operatorname{ran}d$, which is therefore closed. Stokes and
Hilbert complex theory then give the orthogonal decomposition
$L^2=\cH\oplus\operatorname{ran}d\oplus\operatorname{ran}\delta$
and uniqueness of harmonic representatives.

\emph{The action of ordinary cohomology.}
An abstract isomorphism of vector spaces would not yet identify
wedge product by a K\"ahler form with its Lefschetz operator on
$\IH$. For that purpose we use a resolution of the constant
sheaf whose sections can multiply the $L^2$ complex.
For $\mathcal A^\bullet$, normality makes sufficiently small
regular loci connected, so the degree-zero kernel is $\RR$.
In a product-cone chart, simultaneous radial contraction $m_t$
in all factors gives, for a bounded closed $j$-form with $j>0$,
\[
 K\alpha=\int_0^1m_t^*(\iota_{r\partial_r}\alpha)\,\frac{dt}{t},
 \qquad |K\alpha|\le\frac rj\|\alpha\|_\infty.
\]
The distributional homotopy identity gives $dK\alpha=\alpha$;
$m_t^*\alpha\to0$ since $j>0$. Quasi-isometry transfers this
to the actual chart. Thus $\mathcal A^\bullet$ is a fine
resolution. Wedge product with a bounded form and bounded weak
derivative preserves $\mathcal L^\bullet$ and satisfies the
Leibniz rule by local regularisation on the regular locus.
The resulting pairing extends scalar multiplication by $\RR$.
Over a field, tensoring the quasi-isomorphism
$\RR_M\to\mathcal A^\bullet$ with
$\mathcal L^\bullet$ preserves quasi-isomorphisms. The composite
\[
 \RR_M\otimes\mathcal L^\bullet
 \longrightarrow\mathcal A^\bullet\otimes\mathcal L^\bullet
 \longrightarrow\mathcal L^\bullet
\]
is the identity scalar action. The terms of these resolutions are
fine, so their hypercohomology products are computed on global
sections. Wedge product therefore gives precisely the canonical
$H^\bullet(M;\RR)$ module action \cite[II.7]{bredon97}.
Naturality under $\mathcal L^\bullet\simeq\IC_M$ proves
the asserted compatibility.
\end{proof}

\subsection{The pairing}
The module identification concerns multiplication by bounded
representatives. Two arbitrary $L^2$ forms need not have an
$L^2$ wedge product, but their product in complementary degrees
is integrable. Integration thus gives the duality pairing by a
different argument. We verify its normalisation on the regular
locus before invoking the uniqueness of intersection duality.

Put $m=\dim_{\RR}M=16k+4$ and let $\mathbb D_M$ denote the
dualising complex. We use the unsigned intersection pairing
$Q:\IH^j\times\IH^{m-j}\to\RR$. Poincar\'e duality for
$L^2$ cohomology by integration is classical for Witt spaces
with iterated conical metrics \cite{cheeger83b,almp18}.
For the present metric the identification with the intersection
pairing can be proved directly from Stokes and the uniqueness
of the self-duality of the intersection complex.

\begin{proposition}\label{prop:pairing}
If closed $\alpha,\beta\in\Dom d_{\max}$ represent
$\Phi(a),\Phi(b)$ in complementary degrees, then
\begin{equation}\label{eq:pairing}
 Q(a,b)=\int_{M_{\reg}}\alpha\wedge\beta.
\end{equation}
\end{proposition}
\begin{proof}
The product is integrable by Cauchy--Schwarz. Stokes gives, for
$u,v\in\Dom d_{\max}$ with $\deg u=p$ and
$\deg u+\deg v=m-1$,
\[
 \int du\wedge v+(-1)^p\int u\wedge dv=0.
\]
Indeed approximate $u$ in its graph norm by smooth forms
compactly supported in the regular locus and use the weak
derivative of $v$. Thus integration depends only on classes.

For completeness this functional determines the correct
sheaf-theoretic pairing, not just an abstract nondegenerate
form. Over $\RR$, the tensor product
$\mathcal F=\mathcal L^\bullet\otimes\mathcal L^\bullet$
represents $\IC_M\otimes^L\IC_M$ and has fine terms.
Compact Verdier duality gives
\[
 \Hom_{D(M)}(\mathcal F,\mathbb D_M[-m])
 \cong\Hom_{\RR}(H^m\Gamma(M,\mathcal F),\RR).
\]
Integration of the wedge product defines the right-hand
functional. To verify that it kills coboundaries, use a finite
partition of unity to express a section as sums of tensors
of global sections and apply the displayed Stokes identity.
On an open regular subset, the induced functional on compact
supports is ordinary integration; hence its morphism restricts
to the orientation pairing. Self-duality of $\IC_M$ identifies
the space of such pairings with $\Hom(\IC_M,\IC_M)$.
An endomorphism of the intersection complex is determined by
its restriction to the dense regular locus
\cite[\S\S3.5 and 5.3]{gm83}. The morphism is therefore the
intersection pairing. Evaluating on $\alpha\otimes\beta$
proves \eqref{eq:pairing}.
\end{proof}

In top degree, integration is consequently an isomorphism
$H^m_{(2)}\to\RR$. It also follows directly from Stokes,
$\dim\IH^m=1$ and the positive integral of the top K\"ahler
power. For real harmonic forms the Hodge star preserves
$\cH$, and
\begin{equation}\label{eq:positive-pairing}
 \langle\alpha,\beta\rangle_{L^2}=Q(\alpha,*\beta).
\end{equation}
Moreover $*$ belongs on $\cH$ to the associative algebra
generated by $L_1,\Lambda_1$ and $W_1$. To see this, the
Lefschetz decomposition separates the summands
$L_1^r\cH^j_{\rm prim}$. Their projections belong to the
algebra generated by the $\mathfrak{sl}_2$ action, either by
its Casimir and $h$ or by the double-centraliser theorem.
Weil's star formula sends $L_1^r\eta$ to a scalar multiple of
$L_1^{2n-j-r}W_1\eta$ on each such summand
\cite[Proposition~1.2.31]{huybrechts05}. Finally $W_1$ or its
inverse is $\exp(\pi K_{23}/2)$ by the pointwise relations.
Thus an endomorphism commuting with the LLV algebra commutes
with $*$ on $\cH$.

\subsection{The positive three-plane}
The forms $\omega_a$ are parallel and bounded. They define
classes $x_a\in H^2(M;\RR)$ through the resolution
$\mathcal A^\bullet$. The restriction to $M_{\reg}$ is
injective on $H^2$: terminal factoriality gives
$H^2(M)\cong\IH^2(M)$, and codimension at least three gives
$\IH^2(M)\cong H^2(M_{\reg})$
\cite[Lemma~2.10, Proposition~2.17]{tighe22}.
Recollection \ref{thm:imported} therefore identifies $x_1$
as a positive multiple of an ample class and $x_2+ix_3$ as
a nonzero multiple of the symplectic class.

\begin{proposition}\label{prop:qF}
There is $c>0$ such that $q_M(x_a,x_b)=c\delta_{ab}$.
\end{proposition}
\begin{proof}
The module statement and top-degree integration show that
evaluation of a product of bounded closed representatives on
$[M]$ is a fixed positive multiple of their wedge integral.
On a quaternionic tangent space, unit quaternions act by
volume-preserving isometries rotating the three K\"ahler
forms. Thus
\[
 \int(t_1\omega_1+t_2\omega_2+t_3\omega_3)^{2n}
 =C(t_1^2+t_2^2+t_3^2)^n,\qquad C>0.
\]
The Fujiki relation \cite[Proposition~5.15]{bl22} makes this
a positive multiple of $q_M(\sum t_ax_a)^n$. This quadratic
form is nonzero away from zero on the three-plane and positive
at the ample ray. It is therefore the positive $n$th root
of the displayed polynomial, giving the claim by polarisation.
\end{proof}

Theorem \ref{thm:T1}, Proposition \ref{prop:pairing} and
Corollary \ref{thm:T2} prove all of Theorem
\ref{thm:analytic}. Section \ref{sec:algebra} now proves
Theorem \ref{thm:main} without any remaining analytic premise.
